\section{Preguntas abiertas del ecosistema de herramientas de desarrollo}

\subsection{Gran integración o verticalización}

\begin{itemize}
    \item ¿Se fusionarán en una sola plataforma integral las soluciones puntuales de revisión de código, construcción de aplicaciones, monitoreo, etc.?
    \item ¿Evolucionará el AI IDE hacia una AI terminal, y de ahí hacia un AI browser / Replit?
    \item ¿Se dirigirán herramientas como Warp y Cursor hacia la verticalización, enfocándose en dominios específicos como el desarrollo full-stack?
\end{itemize}

\subsection{Estándares para la personalización de agentes}

\begin{importantbox}{De la fragmentación a un estándar unificado}
El panorama actual, fragmentado y disperso entre .cursorrules, CLAUDE.md, agents.md y otros, no durará mucho tiempo. AGENTS.md es el primer paso hacia un estándar unificado. En el futuro surgirá un estándar común para personalizar al coding agent y proporcionar reglas específicas de cada proyecto.
\end{importantbox}

\subsection{La doble disputa entre la capa de producto y la capa de infraestructura}

Un punto importante del curso, y que suele pasarse por alto, es que las herramientas de desarrollo con IA no son un problema puramente de producto en el frontend; también están limitadas por el modelo subyacente y por la infraestructura de inferencia.

\begin{table}[H]
\centering
\begin{tabular}{p{0.23\textwidth} p{0.31\textwidth} p{0.34\textwidth}}
\toprule
Dimensión & Decisión a tomar & Impacto directo en el producto \\
\midrule
Selección del modelo & ¿Usar un frontier model de código cerrado o un modelo de código abierto? & Determina el límite superior de calidad, la estructura de costos y el grado de control \\
Estrategia de contexto & ¿Indexación total o recuperación bajo demanda? & Determina la latencia de respuesta, la precisión de las referencias y el desempeño de la memoria de largo plazo \\
Entorno de ejecución aislado & ¿Ejecución local o contenedor remoto? & Determina la seguridad, los límites de permisos y la profundidad de integración \\
Ecosistema de herramientas & ¿Herramientas integradas o extensiones vía MCP? & Determina la apertura del sistema y el costo de mantenimiento \\
Ciclo de retroalimentación de datos & ¿Se registran la retroalimentación y las trayectorias de edición? & Determina la velocidad de iteración del producto y la capacidad de personalización \\
\bottomrule
\end{tabular}
\caption{Acoplamiento entre la capa de producto y la capa de infraestructura de las herramientas de desarrollo con IA}
\end{table}

\begin{warningbox}{No reduzcas la ``AI UX'' a una ventana de chat}
Muchos equipos creen erróneamente que basta con agregar un chat panel para lograr la conversión hacia IA. Pero la verdadera dificultad está en el ensamblaje del contexto, la invocación de herramientas, el control de permisos, la recuperación de errores y la explicación de los resultados. Si falta esa infraestructura, ni la interfaz más atractiva pasará de ofrecer una experiencia de nivel demostrativo.
\end{warningbox}

\subsection{Marco de evaluación para la adquisición en equipo y la adopción individual}

De la discusión del curso se puede extraer un marco de evaluación más práctico, que ayuda a un equipo a decidir si vale la pena implementar una herramienta de desarrollo con IA:

\begin{enumerate}
    \item \textbf{Time-to-value}: ¿cuánto tiempo pasa, después del primer uso, antes de ahorrar realmente tiempo de trabajo?
    \item \textbf{Trust surface}: cuando el modelo se equivoca, ¿puede el usuario ubicar rápidamente la causa?
    \item \textbf{Workflow fit}: ¿se integra al flujo de trabajo existente, o obliga al equipo a cambiar su manera de trabajar?
    \item \textbf{Admin control}: ¿puede la organización controlar el modelo, los permisos, la auditoría y la retención de datos?
    \item \textbf{Extensibility}: ¿será fácil, en el futuro, conectar herramientas personalizadas, APIs internas y bases de conocimiento?
\end{enumerate}

\begin{knowledgebox}{Por qué a las empresas les importa tanto el ``flujo de trabajo interpretable''}
Los usuarios de consumo pueden tolerar cierta dosis de sorpresa y fallas ocasionales, pero al momento de adquirir una herramienta, a las empresas les preocupa más la cadena de responsabilidad. Una herramienta capaz de explicar cada paso de su actuación, de mostrar el origen del contexto y de reproducir el proceso de ejecución suele ser más fácil de implementar que una herramienta que es más inteligente en un solo uso pero inexplicable.
\end{knowledgebox}

\subsection{De herramienta de productividad individual a estándar de equipo}

El curso también deja implícita una ruta de expansión de producto: muchas herramientas de desarrollo con IA entran primero como herramientas de productividad individual, y solo después intentan convertirse en un estándar de equipo. Los enfoques de ambas etapas no son los mismos.

\begin{itemize}
    \item En la etapa de adopción individual, lo que más importa es el momento de asombro (wow moment) y el beneficio inmediato
    \item En la etapa de despliegue en equipo, importan más los permisos, la auditoría, la configurabilidad y la consistencia
    \item Cuando una herramienta empieza a formar parte del flujo estándar de un equipo, debe ofrecer reglas compartibles, registros de ejecución reproducibles y vistas de administración
\end{itemize}

\begin{warningbox}{Que la experiencia de un solo usuario sea excelente no significa que sea viable en la organización}
Muchos productos de IA son muy sólidos en una computadora personal, pero en cuanto entran a la red de una empresa, a su sistema de permisos, a la cadena de revisión de código y a los requisitos de cumplimiento, se exponen sus deficiencias. El verdadero foso competitivo suele aparecer en la capacidad de gobernanza a nivel de equipo, y no en la calidad de una sola generación.
\end{warningbox}

\subsection{Resumen de la sección}

El ecosistema de herramientas de desarrollo está en una etapa de rápida evolución; entre las preguntas clave aún sin resolver están: integración frente a verticalización, la unificación del estándar de personalización de agentes, y los límites de viabilidad del lenguaje natural como interfaz principal de programación.
